\documentclass[11pt]{article}
% The preamble every project's paper shares: % The preamble every project's paper shares: % The preamble every project's paper shares: \input{../../preamble} from
% projects/<project>/paper/. Boards come from chessboard and xskak, plots from
% pgfplots; every number a paper prints comes from its generated.tex.
\usepackage[LSBC4,T1]{fontenc}
\usepackage[utf8]{inputenc}
\usepackage{lmodern}
\usepackage[margin=1in]{geometry}
\usepackage{microtype}
\usepackage{amsmath}
\usepackage{booktabs}
\usepackage{tabularx}
\usepackage{array}
\usepackage{xcolor}
\usepackage[obeyspaces,spaces]{url}
\usepackage{xskak}
\usepackage{pgfplots}
\pgfplotsset{compat=1.18}
\usepackage{caption}
\usepackage{subcaption}
\usepackage{enumitem}
\usepackage[section]{placeins}
\usepackage[hidelinks]{hyperref}

\setchessboard{boardfontencoding=LSBC4, boardfontfamily=skaknew, blackfieldcolor=black!22,
  whitefieldcolor=white, setfontcolors, showmover=false, label=true, labelfontsize=6pt,
  boardfontsize=14pt}

% A move or line in UCI, and a position in FEN, both breakable at a slash.
\newcommand{\uci}[1]{\texttt{#1}}
\newcommand{\fen}[1]{{\footnotesize\path{#1}}}
\newcolumntype{L}{>{\raggedright\arraybackslash}X}

\setlength{\parskip}{0.4em}
\setlength{\parindent}{0pt}
 from
% projects/<project>/paper/. Boards come from chessboard and xskak, plots from
% pgfplots; every number a paper prints comes from its generated.tex.
\usepackage[LSBC4,T1]{fontenc}
\usepackage[utf8]{inputenc}
\usepackage{lmodern}
\usepackage[margin=1in]{geometry}
\usepackage{microtype}
\usepackage{amsmath}
\usepackage{booktabs}
\usepackage{tabularx}
\usepackage{array}
\usepackage{xcolor}
\usepackage[obeyspaces,spaces]{url}
\usepackage{xskak}
\usepackage{pgfplots}
\pgfplotsset{compat=1.18}
\usepackage{caption}
\usepackage{subcaption}
\usepackage{enumitem}
\usepackage[section]{placeins}
\usepackage[hidelinks]{hyperref}

\setchessboard{boardfontencoding=LSBC4, boardfontfamily=skaknew, blackfieldcolor=black!22,
  whitefieldcolor=white, setfontcolors, showmover=false, label=true, labelfontsize=6pt,
  boardfontsize=14pt}

% A move or line in UCI, and a position in FEN, both breakable at a slash.
\newcommand{\uci}[1]{\texttt{#1}}
\newcommand{\fen}[1]{{\footnotesize\path{#1}}}
\newcolumntype{L}{>{\raggedright\arraybackslash}X}

\setlength{\parskip}{0.4em}
\setlength{\parindent}{0pt}
 from
% projects/<project>/paper/. Boards come from chessboard and xskak, plots from
% pgfplots; every number a paper prints comes from its generated.tex.
\usepackage[LSBC4,T1]{fontenc}
\usepackage[utf8]{inputenc}
\usepackage{lmodern}
\usepackage[margin=1in]{geometry}
\usepackage{microtype}
\usepackage{amsmath}
\usepackage{booktabs}
\usepackage{tabularx}
\usepackage{array}
\usepackage{xcolor}
\usepackage[obeyspaces,spaces]{url}
\usepackage{xskak}
\usepackage{pgfplots}
\pgfplotsset{compat=1.18}
\usepackage{caption}
\usepackage{subcaption}
\usepackage{enumitem}
\usepackage[section]{placeins}
\usepackage[hidelinks]{hyperref}

\setchessboard{boardfontencoding=LSBC4, boardfontfamily=skaknew, blackfieldcolor=black!22,
  whitefieldcolor=white, setfontcolors, showmover=false, label=true, labelfontsize=6pt,
  boardfontsize=14pt}

% A move or line in UCI, and a position in FEN, both breakable at a slash.
\newcommand{\uci}[1]{\texttt{#1}}
\newcommand{\fen}[1]{{\footnotesize\path{#1}}}
\newcolumntype{L}{>{\raggedright\arraybackslash}X}

\setlength{\parskip}{0.4em}
\setlength{\parindent}{0pt}

% Written by projects/phone-vs-routine/paper/generate.mjs; do not edit by hand.

\newcommand{\AppCommit}{d8b771238c}
\newcommand{\Engine}{stockfish 18.0.8 stockfish-18-lite-single}
\newcommand{\Cases}{60}
\newcommand{\Line}{5}
\newcommand{\Draws}{2,000}
\newcommand{\Seed}{20261014}
\newcommand{\PassShare}{25}
\newcommand{\DesktopKnps}{500}
\newcommand{\AcceptWindow}{3}
\newcommand{\PhoneSweep}{100,000}
\newcommand{\DesktopSweep}{250,000}
\newcommand{\RoutineSweep}{1 million}
\newcommand{\VerifyNodes}{750,000}
\newcommand{\VerdictTable}{%
\begin{tabular}{lrrrrlll}
\toprule
 & \multicolumn{2}{c}{serious errors} & & & & & \\
\cmidrule(lr){2-3}
profile & routine & device & missed & extra & miss share & extra share & verdict \\
\midrule
Phone & 341 & 336 & 31 & 26 & 9.1\% {\footnotesize (5.9, 12.6)} & 7.7\% {\footnotesize (4.9, 10.7)} & differs \\
Desktop & 341 & 346 & 25 & 30 & 7.3\% {\footnotesize (4.4, 10.5)} & 8.7\% {\footnotesize (5.7, 11.7)} & differs \\
\bottomrule
\end{tabular}
}
\newcommand{\OtherTable}{%
\begin{tabular}{llllll}
\toprule
 & \multicolumn{2}{c}{blunders} & \multicolumn{2}{c}{inaccuracies} & any class \\
\cmidrule(lr){2-3}\cmidrule(lr){4-5}
profile & missed & extra & missed & extra & differs \\
\midrule
Phone & 8.1\% {\footnotesize (3.7, 12.3)} & 6.0\% {\footnotesize (2.2, 10.0)} & 35.8\% {\footnotesize (29.4, 42.5)} & 27.0\% {\footnotesize (21.9, 32.0)} & 26.1\% {\footnotesize (24.0, 28.2)} \\
Desktop & 6.6\% {\footnotesize (2.7, 10.7)} & 6.6\% {\footnotesize (2.8, 10.8)} & 32.7\% {\footnotesize (27.0, 38.8)} & 24.6\% {\footnotesize (19.2, 29.9)} & 23.1\% {\footnotesize (21.1, 25.2)} \\
\bottomrule
\end{tabular}
}
\newcommand{\WhereTable}{%
\begin{tabular}{lrrrrrr}
\toprule
 & \multicolumn{3}{c}{missed} & \multicolumn{3}{c}{extra} \\
\cmidrule(lr){2-4}\cmidrule(lr){5-7}
profile & second look & sweep only & routine: tablebase & second look & sweep only & routine: tablebase \\
\midrule
Phone & 26 & 5 & 0 & 26 & 0 & 1 \\
Desktop & 23 & 2 & 0 & 30 & 0 & 0 \\
\bottomrule
\end{tabular}
}
\newcommand{\PhoneVerdict}{differs}
\newcommand{\PhoneGames}{60}
\newcommand{\PhonePlies}{3,917}
\newcommand{\PhoneRoutineSerious}{341}
\newcommand{\PhoneDeviceSerious}{336}
\newcommand{\PhoneBoth}{310}
\newcommand{\PhoneMissed}{31}
\newcommand{\PhoneExtra}{26}
\newcommand{\PhoneMiss}{9.1}
\newcommand{\PhoneMissLo}{5.9}
\newcommand{\PhoneMissHi}{12.6}
\newcommand{\PhoneExtraShare}{7.7}
\newcommand{\PhoneExtraLo}{4.9}
\newcommand{\PhoneExtraHi}{10.7}
\newcommand{\PhoneDiffer}{26.1}
\newcommand{\PhoneDifferLo}{24.0}
\newcommand{\PhoneDifferHi}{28.2}
\newcommand{\PhoneInaccMiss}{36}
\newcommand{\PhoneInaccExtra}{27}
\newcommand{\PhoneLooks}{999}
\newcommand{\PhoneFlags}{1,788}
\newcommand{\PhoneCut}{789}
\newcommand{\PhoneBias}{$-$0.6}
\newcommand{\PhoneMae}{2.0}
\newcommand{\PhoneLargest}{6.4}
\newcommand{\PhoneSides}{120}
\newcommand{\PhoneMissedSweep}{five}
\newcommand{\PhoneMissedVerified}{26}
\newcommand{\PhoneExtraSweep}{zero}
\newcommand{\PhoneExtraVerified}{26}
\newcommand{\PhoneMissedTablebase}{zero}
\newcommand{\PhoneExtraTablebase}{one}
\newcommand{\DesktopVerdict}{differs}
\newcommand{\DesktopGames}{60}
\newcommand{\DesktopPlies}{3,917}
\newcommand{\DesktopRoutineSerious}{341}
\newcommand{\DesktopDeviceSerious}{346}
\newcommand{\DesktopBoth}{316}
\newcommand{\DesktopMissed}{25}
\newcommand{\DesktopExtra}{30}
\newcommand{\DesktopMiss}{7.3}
\newcommand{\DesktopMissLo}{4.4}
\newcommand{\DesktopMissHi}{10.5}
\newcommand{\DesktopExtraShare}{8.7}
\newcommand{\DesktopExtraLo}{5.7}
\newcommand{\DesktopExtraHi}{11.7}
\newcommand{\DesktopDiffer}{23.1}
\newcommand{\DesktopDifferLo}{21.1}
\newcommand{\DesktopDifferHi}{25.2}
\newcommand{\DesktopInaccMiss}{33}
\newcommand{\DesktopInaccExtra}{25}
\newcommand{\DesktopLooks}{994}
\newcommand{\DesktopFlags}{1,786}
\newcommand{\DesktopCut}{792}
\newcommand{\DesktopBias}{$-$0.7}
\newcommand{\DesktopMae}{1.6}
\newcommand{\DesktopLargest}{7.3}
\newcommand{\DesktopSides}{120}
\newcommand{\DesktopMissedSweep}{two}
\newcommand{\DesktopMissedVerified}{23}
\newcommand{\DesktopExtraSweep}{zero}
\newcommand{\DesktopExtraVerified}{30}
\newcommand{\DesktopMissedTablebase}{zero}
\newcommand{\DesktopExtraTablebase}{zero}
\newcommand{\RoutineLooks}{1,738}
\newcommand{\PhoneBiasAbs}{0.6}
\newcommand{\DesktopBiasAbs}{0.7}
\newcommand{\PhoneOneIn}{eleven}
\newcommand{\Near}{3}
\newcommand{\MistakeLine}{10}
\newcommand{\ReadingsTable}{%
\begin{tabular}{lrrrrrrrr}
\toprule
 & \multicolumn{4}{c}{moves both looked at twice} & \multicolumn{4}{c}{serious errors they disagree on} \\
\cmidrule(lr){2-5}\cmidrule(lr){6-9}
 & & same & same & loss gap & & mistake & both within & an extra look \\
profile & moves & depth & numbers & median (90\%) & moves & /inaccuracy & 3 points & changed it \\
\midrule
Phone & 880 & 260 & 17 & 0.57 (2.6) & 57 & 50 & 40 & 26 \\
Desktop & 905 & 245 & 18 & 0.63 (2.4) & 55 & 54 & 49 & 25 \\
\bottomrule
\end{tabular}
}
\newcommand{\PhoneBothLooked}{880}
\newcommand{\PhoneSameDepth}{260}
\newcommand{\PhoneIdentical}{17}
\newcommand{\PhoneGapMedian}{0.6}
\newcommand{\PhoneGapNinety}{2.6}
\newcommand{\PhoneDisagree}{57}
\newcommand{\PhoneFlips}{50}
\newcommand{\PhoneFlipShare}{88}
\newcommand{\PhoneNear}{40}
\newcommand{\PhoneUnsettled}{26}
\newcommand{\DesktopBothLooked}{905}
\newcommand{\DesktopSameDepth}{245}
\newcommand{\DesktopIdentical}{18}
\newcommand{\DesktopGapMedian}{0.6}
\newcommand{\DesktopGapNinety}{2.4}
\newcommand{\DesktopDisagree}{55}
\newcommand{\DesktopFlips}{54}
\newcommand{\DesktopFlipShare}{98}
\newcommand{\DesktopNear}{49}
\newcommand{\DesktopUnsettled}{25}
\newcommand{\PvDMiss}{9.8}
\newcommand{\PvDMissLo}{6.5}
\newcommand{\PvDMissHi}{13.4}
\newcommand{\PvDExtra}{7.1}
\newcommand{\PvDExtraLo}{4.6}
\newcommand{\PvDExtraHi}{9.9}
\newcommand{\VerdictLine}{314}
\newcommand{\VerdictFile}{\texttt{src/services/livecoach/verdict.ts}}
\newcommand{\GameReviewFile}{\texttt{src/features/gamereview/GameReview.tsx}}
\newcommand{\KeptLine}{203}
\newcommand{\CoachLine}{214}
\newcommand{\DeviceHash}{64}
\newcommand{\NodeHash}{128}
\newcommand{\DeviceEngineFile}{\texttt{src/adapters/browser/engine/index.ts}}
\newcommand{\ExampleBoard}{%
\chessboard[setfen={6k1/5ppp/P3p1b1/3p2q1/3rP3/5P2/4Q1PP/R5K1 b - - 0 28}, boardfontsize=13pt, pgfstyle=straightmove, arrow=to, linewidth=0.07em, color=red!75!black, markmoves={g5-d8}, pgfstyle=straightmove, arrow=to, linewidth=0.07em, color=green!50!black, markmoves={d4-c4}]
}
\newcommand{\ExampleFen}{\fen{6k1/5ppp/P3p1b1/3p2q1/3rP3/5P2/4Q1PP/R5K1 b - - 0 28}}
\newcommand{\ExampleMove}{28\ldots{}\symqueen{}d8}
\newcommand{\ExampleUci}{\uci{g5d8}}
\newcommand{\ExampleBest}{28\ldots{}\symrook{}c4 29.~\symqueen{}xc4 dxc4 30.~a7}
\newcommand{\ExampleBestUci}{\uci{d4c4 e2c4 d5c4 a6a7}}
\newcommand{\ExampleRoutineCls}{blunder}
\newcommand{\ExamplePhoneCls}{inaccuracy}
\newcommand{\ExampleRoutineLoss}{22.3}
\newcommand{\ExamplePhoneLoss}{6.2}
\newcommand{\ExampleSide}{Black}
\newcommand{\ExampleDepth}{16}
\newcommand{\ExampleBestMove}{28\ldots{}\symrook{}c4}
\newcommand{\ExamplePhonePlayed}{$+$1.68}
\newcommand{\ExampleRoutinePlayed}{$+$3.62}
\newcommand{\ExampleCandidates}{three}


\title{Does Game Review on his phone agree with the routine?\\[0.3em]
\large The phone's, the desktop's and the routine's reviews of the same \Cases{} games}
\author{Chess Coach research lab}
\date{}

\begin{document}
\maketitle

\begin{abstract}
On \Cases{} of his games, Game Review on a phone or a desktop does not name the same serious errors
as the routine's review of the same game. By the rule committed before any comparison, both devices
differ from the routine: the phone misses \PhoneMiss\% of the routine's mistakes and blunders and
names \PhoneExtraShare\% of its own that the routine does not; the desktop misses \DesktopMiss\%
and adds \DesktopExtraShare\%. The next step the rule named, a better choice of which moves get the
deeper second look, would not close the gap, because the devices looked twice at nearly every move
they disagree on. In checks added after the results were read, nearly every disagreement is a
mistake on one side and an inaccuracy on the other, and the second look itself does not repeat: on
the \PhoneBothLooked{} moves both the phone and the routine looked at twice, with the same nodes from
the same position, the two looks gave the same numbers on \PhoneIdentical. The phone and the desktop
disagree with each other as much as either does with the routine. What tells the reviews apart is
what the engine searched earlier in the game, not the phone's smaller budget. A fact about the code
makes the difference last: Game Review keeps the review a device made of a game, and opens the
routine's only when it has kept none. Opening the routine's once there is one would end the
difference for every game the routine reviews; a second look that starts with the engine's memory
cleared is the next round.
\end{abstract}

\section{The question}

The app reviews a game in two places. On a device, Game Review sweeps each position at
\PhoneSweep{} nodes on a phone or \DesktopSweep{} on a desktop (a device counts as a desktop at
\DesktopKnps{} thousand nodes a second or more), gives the deeper second look to at most
\PassShare\% of the game's moves, and reads no tablebase. The daily routine sweeps at
\RoutineSweep{} nodes, gives every move it flags the second look and reads the Lichess tablebase.
The second look itself is \VerifyNodes{} nodes everywhere. When he opens a game on his phone, does
Game Review name the same mistakes and blunders as the routine's review of it? Where it does not, is
the move one the device never looked at twice?

\begin{table}[h]
\centering
\small
\begin{tabular}{llll}
\toprule
profile & sweep (nodes) & second look & tablebase \\
\midrule
phone & \PhoneSweep & at most \PassShare\% of moves, \VerifyNodes{} nodes & no \\
desktop & \DesktopSweep & at most \PassShare\% of moves, \VerifyNodes{} nodes & no \\
routine & \RoutineSweep & every flagged move, \VerifyNodes{} nodes & yes \\
\bottomrule
\end{tabular}
\caption{The three review profiles, as the app at commit \texttt{\AppCommit} sets them.}
\label{tab:profiles}
\end{table}

\section{The data}

The \Cases{} games of the research records' frozen sample, drawn from his archive by length and
result. Each was reviewed by the app's own entry point (commit \texttt{\AppCommit}, \Engine) at the
phone, desktop and routine profiles, with a fresh engine and no cache for every game, so a document
depends only on the code and the game; on one thread at fixed nodes the engine is deterministic.

\section{The method}

All the measures come from the review documents, with the app's own definitions. A ply is
\emph{serious} when the metrics service's \texttt{isSerious} says so of its class: a mistake, a
blunder or a miss. A serious error the routine names is \emph{missed} by a device when the device's
record of the same ply is not serious; a device's serious error is \emph{extra} when the routine's
record of that ply is not. The \emph{miss share} is the missed plies over the routine's serious
errors, the \emph{extra share} the extra plies over the device's. Beside them: the same shares for
blunders and for inaccuracies, the plies whose class differs in any way, and each side's accuracy on
the device against the routine. For every missed or extra ply, the device either gave it the second
look or left it at the sweep's reading, and the routine's reading either came from the tablebase or
did not.

Shares are pooled over the \Cases{} games, with 95\% intervals from \Draws{} bootstrap draws over
games (a Poisson(1) weight for each game, seed \Seed).

The rule, committed before any comparison was made, for each device: the device \emph{agrees} with
the routine when the upper ends of both the miss share's and the extra share's intervals lie below
\Line\%; it \emph{differs} when the lower end of either lies above \Line\%; otherwise the games
cannot tell (\emph{unclear}). One in twenty is the line: a phone that missed more would, every few
games, leave out an error the daily review names, which a player reading both would notice within a
week. When a device differs and most of its missed or extra plies are ones it left at the sweep's
reading, the next round chooses better which moves get the second look, at the device's fixed
budget.

Four descriptions were added after the results were read, and each is marked where it appears:
whether the two profiles' second looks of the same move give the same numbers, how far apart their
losses are, how near the mistake line the disagreements lie, and the phone against the desktop. They
test nothing and carry no rule.

\section{Results}

\subsection{The verdict}

By the rule, both devices differ from the routine (Table~\ref{tab:verdict}). The phone misses
\PhoneMissed{} of the routine's \PhoneRoutineSerious{} serious errors, \PhoneMiss\%
(\PhoneMissLo{} to \PhoneMissHi\%), and names \PhoneExtra{} that the routine does not,
\PhoneExtraShare\% of its own (\PhoneExtraLo{} to \PhoneExtraHi\%). The desktop misses
\DesktopMiss\% (\DesktopMissLo{} to \DesktopMissHi\%) and adds \DesktopExtraShare\%
(\DesktopExtraLo{} to \DesktopExtraHi\%). The phone's miss share and the desktop's extra share each
lie wholly above the \Line\% line. Blunders alone part about as often, and about a third of the
inaccuracies differ (Table~\ref{tab:other}): on the phone, \PhoneDiffer\% of all moves get a
different class from the routine's. Both devices read a side's accuracy a little lower than the
routine does, by \PhoneBiasAbs{} and \DesktopBiasAbs{} points on average.

\begin{table}[h]
\centering
\scriptsize\VerdictTable
\caption{The rule: serious errors (mistakes, blunders and misses) named by the routine and by each
device over the \Cases{} games, the routine's that the device missed, the device's that the routine
does not name, the two shares with their 95\% intervals, and the verdict against the \Line\% line.}
\label{tab:verdict}
\end{table}

\begin{table}[h]
\centering
\scriptsize\OtherTable
\caption{Beside the rule: the same shares for blunders alone and for inaccuracies, and the share of
all moves whose class differs in any way, with 95\% intervals.}
\label{tab:other}
\end{table}

\subsection{Where the differences come from}

Not from the second look's quota. The phone gave the second look to \PhoneMissedVerified{} of the
\PhoneMissed{} serious errors it missed and to all \PhoneExtra{} it added; the desktop to
\DesktopMissedVerified{} of \DesktopMissed{} and all \DesktopExtra{} (Table~\ref{tab:where}). The
tablebase played almost no part. So the next round the rule named, choosing better which moves get
the second look, would leave nearly every difference where it is.

\begin{table}[h]
\centering
\scriptsize\WhereTable
\caption{Where each missed or extra serious error comes from: whether the device gave the move the
second look or left it at the sweep's reading, and whether the routine's reading came from the
tablebase.}
\label{tab:where}
\end{table}

\subsection{Why the reviews part}

The second look is the same search on every profile: \VerifyNodes{} nodes and the same candidates,
from the same position with the game's history. But the engine keeps what it learned from the game's
earlier searches, and each profile has searched different things before it reaches a move: the
sweep at its own budget, and the second looks it chose before. In a check added after the results
were read, on the \PhoneBothLooked{} moves both the phone and the routine looked at twice, the two
looks gave the same numbers on only \PhoneIdentical{} and reached the same depth on
\PhoneSameDepth{} (Table~\ref{tab:readings}). Their losses for the move differ by \PhoneGapMedian{}
points at the median and \PhoneGapNinety{} at the 90th percentile.

That is enough to move a class at the line. Of the \PhoneDisagree{} serious errors the phone and the
routine disagree on, \PhoneFlips{} are a mistake on one side and an inaccuracy on the other, and
\PhoneNear{} have both losses within \Near{} points of the \MistakeLine-point mistake line; on
\PhoneUnsettled, one profile's own extra look, at twice the nodes, changed its verdict, so the engine
had not settled them either. The desktop's are alike. And the phone and the desktop, whose profiles
differ only in the sweep, disagree with each other as much as either does with the routine: the
phone misses \PvDMiss\% of the desktop's serious errors (\PvDMissLo{} to \PvDMissHi\%) and adds
\PvDExtra\% (\PvDExtraLo{} to \PvDExtraHi\%).

\begin{table}[h]
\centering
\scriptsize\ReadingsTable
\caption{Why the reviews part (added after the results were read; it tests nothing). On the moves
both the device and the routine looked at twice: how many, on how many the two looks reached the
same depth or gave the same numbers (the best move's score and the played move's), and the gap
between their losses. On the serious errors they disagree on: how many, how many are a mistake on one
side and an inaccuracy on the other, how many have both losses within \Near{} points of the
\MistakeLine-point line, and on how many either profile's extra look changed its verdict.}
\label{tab:readings}
\end{table}

\subsection{An example}

The phone's costliest miss is \ExampleSide's \ExampleMove{} in Figure~\ref{fig:example}. The
routine calls it a \ExampleRoutineCls{} (a loss of \ExampleRoutineLoss{} points) and the phone an
\ExamplePhoneCls{} (\ExamplePhoneLoss). Both looked at it twice, and both looks found
\ExampleBestMove{} best at depth \ExampleDepth. They part on the move played. The phone's look kept
\ExampleMove{} among its \ExampleCandidates{} candidates and scored it \ExamplePhonePlayed{} there;
the routine's left it out and searched it on its own, at \ExampleRoutinePlayed{} (both from White's
side, in pawns). The routine's line is \ExampleBest{} (\ExampleBestUci), from the position
\ExampleFen.

\begin{figure}[h]
\centering
\ExampleBoard
\caption{\ExampleSide{} to move: \ExampleMove{} (red) where the routine's second look prefers
\ExampleBestMove{} (green).}
\label{fig:example}
\end{figure}

\section{Facts about the code}

These are facts about the app at commit \texttt{\AppCommit}, not measurements.

\paragraph{Which review Game Review opens.} When he opens a game, Game Review first reads the review
it kept on the device (\GameReviewFile, line \KeptLine) and asks for the coach's own review only
when it has kept none (line \CoachLine). The daily routine reviews the evening's games the next
morning. A game he opens on his phone before then is reviewed on the phone, and the phone keeps that
review: once the routine has reviewed the game, the daily review names the routine's serious errors
and Game Review still shows the phone's. For those games the differences measured here stay. A game
he first opens after the routine reviewed it shows the routine's review.

\paragraph{The device's engine.} On a device the engine runs with a \DeviceHash{} MB hash and keeps
an evaluation cache (\DeviceEngineFile); the routine's runs with a \NodeHash{} MB hash. This study
ran every profile on the routine's engine with no cache, so here the phone and the desktop differ
from the routine only by their profile.

\paragraph{The live coach's accept window.} The Play screen's live coach decides whether the move
played was among the fine ones with its own \AcceptWindow{} (\VerdictFile, line \VerdictLine), where
the review imports \texttt{ACCEPT\_WINDOW\_WP}, also \AcceptWindow. The two agree today, but a change
to the review's window would not reach the live coach.

\section{What it means for the app}

\begin{enumerate}[leftmargin=*]
\item \textbf{The phone's review and the routine's part on about one serious error in
\PhoneOneIn.} Nearly all of them are moves near the line between a mistake and an inaccuracy. A
player reading both would meet \PhoneDisagree{} such moves in \Cases{} games.
\item \textbf{Open the coach's review in Game Review once there is one,} in place of a review the
device kept, so that his phone and the daily review name the same serious errors for every game the
routine has reviewed. Before the routine reaches a game, the device's review is the only one there
is.
\item \textbf{A bigger device budget would not fix it.} The desktop parts from the routine as much
as the phone does, and the phone and the desktop part from each other as much. The second-look quota
is not the cause either.
\item \textbf{Make the second look repeat,} for the games the routine has not reached. With the
engine's memory cleared before each second look, on a hash of one size, the look becomes a function
of the position, the history, the nodes and the engine build, and every profile that looks at a move
twice would read it the same way. The cost is that the look no longer
starts from the sweep's work, so it may search less deeply for its nodes. The next round measures
that change on these \Cases{} games against this paper's rule, with its own rule committed first.
\item \textbf{The live coach imports the review's accept window} instead of typing its own
\AcceptWindow, so that a change to one reaches the other.
\end{enumerate}

\section{Limits}

The \Cases{} games are his games of one stretch of the archive, sampled for an accuracy study. The routine
read here is today's: the daily reviews he has read were written by the routine on the app of their
day, and this compares the profiles on one commit, so a difference is the profiles', not the app's
history. Which profile his phone runs depends on its measured speed, so both device profiles are
reported. Every profile ran on the routine's engine with no cache; a real device's engine has a
\DeviceHash{} MB hash and a cache, so its second looks start from a memory that differs further,
which this study did not measure. The descriptions of why the reviews part came after the results
were read; they say what the measured disagreements look like and test nothing.

\end{document}
